% !TEX root = ../main.tex

\begin{abstract}[zh]
  在云原生与事件驱动架构广泛应用的背景下，Apache Flink 与 Apache Kafka 的组合已成为实时数据流水线的常见实现方式。
  然而，现有 Flink Kafka Sink 多采用静态配置模式，作业启动后目标 Topic 与相关连接信息通常被固化；
  当下游 Kafka 拓扑因扩容、迁移、容灾切换或业务演进而发生变化时，往往需要停机修改配置并重新部署，
  从而增加运维成本并削弱实时链路的连续性。相比之下，社区在消费端已经探索了动态 Kafka Source，
  而生产端动态 Sink 的研究与实现仍相对不足。

  本文围绕“基于 Flink 的动态 Kafka Topic 发现与订阅机制”展开研究。首先，从统一 Sink API、
  Kafka 元数据访问与静态 KafkaSink 的使用方式出发，分析其在动态写入场景下的局限；
  其次，结合开题报告与中期实现，提出基于正则表达式的 Topic 发现方案，以及逻辑路由与物理路由分离的设计思路；
  在此基础上，依托现有代码实现了 DynamicKafkaSink、DynamicKafkaSinkWriter、
  SingleClusterTopicMetadataService、DynamicKafkaSinkWriterState、DynamicKafkaCommittable
  与 DynamicKafkaCommitter 等关键组件，并通过广播状态与 route 更新事件支持运行期动态路由。
  该原型能够在不重启 Flink 作业的前提下发现符合规则的 Topic、建立对应写入路径，并完成多目标写入与初步状态恢复。

  目前，系统已在本地环境完成基本功能验证，验证了动态 Topic 发现、显式 route 更新与多目标写入机制的可行性。
  同时，本文也对当前实现中的局限进行了分析，包括元数据轮询开销、失效 route 的 writer 回收不足、
  Exactly-once 语义验证尚不完整等问题。后续工作将围绕元数据访问优化、资源管理、故障恢复一致性验证
  及与静态 Sink 的对比实验继续推进。
\end{abstract}
